% Exploratory conditional reduction, outside the frozen CP27 paper.
% Root written after its 23:00 UTC source and local-proof checks.
\subsection*{A conditional single-orbit inertia route for odd degrees}

This note does not assert a new monodromy image or a smaller proved main
gate. The required global branch-divisor assertions below are unproved.
The two elementary lemmas are supplied with complete proofs for independent
review. In particular an irreducible branch divisor alone is insufficient.

\paragraph{Algebraic criterion.}
Let $V$ be a nonzero finite-dimensional vector space over a field $k$, and
let $G\subseteq\operatorname{GL}_k(V)$ be generated by a single
$G$-conjugacy class $\mathcal C$ of transvections. If
$V^G=0=(V^\vee)^G$, then $V$ is an irreducible $k[G]$-module.

\emph{Proof.}
Suppose $0\ne S\subsetneq V$ is $G$-stable. If every element of
$\mathcal C$ acts identically on $S$, its generators do so, hence
$S\subseteq V^G=0$, a contradiction. Otherwise choose $t\in\mathcal C$
and $s\in S$ with $(t-1)s\ne0$. Since $t-1$ has rank one,
$\operatorname{im}(t-1)=k(t-1)s\subseteq S$. Conjugacy and stability
then give $\operatorname{im}(u-1)\subseteq S$ for every
$u\in\mathcal C$. Thus every generator acts identically on $V/S$.
Any nonzero linear functional on $V/S$ pulls back to a nonzero element
of $(V^\vee)^G$, again a contradiction. This proves irreducibility.

\paragraph{Local corank-one inertia.}
Let $R$ be a strictly henselian discrete valuation ring of characteristic
two, with fraction field $K$, and let $E$ be a finite free $R$-module
with a Frobenius-semilinear operator $F$. Suppose the solution spaces
of $F(v)=v$ on a geometric generic and special fibre have respective
$\mathbf F_2$-dimensions $r$ and $r-1$. Then the inertia image on the
generic solution space $V$ fixes pointwise a hyperplane $W$ and embeds
in the additive group
$\operatorname{Hom}_{\mathbf F_2}(V/W,W)$. Every nonidentity inertia
element is a transvection. Its image need not be cyclic.

\emph{Proof.}
In a basis, the Lang zero scheme has equations $A x^{[2]}-x=0$.
Their Jacobian is $-I$, so it is etale over $R$, even though it need
not be finite over $R$. Each special solution therefore has a unique
lift through its henselian neighbourhood to an $R$-section. A relation
between the generic lifts reduces to the same relation in the special
fibre; hence the lifts form an $(r-1)$-dimensional subspace $W$.
They are $K$-rational, so inertia fixes $W$ pointwise. On the
one-dimensional quotient $V/W$, inertia acts through
$\mathbf F_2^\times=\{1\}$. Consequently every $g-1$ factors from
$V/W$ into $W$, and these factorizations multiply by addition since
$(g-1)(h-1)=0$. A nonzero factorization has rank one and square zero.
This proves the claims. In particular the permitted image can have
order as large as $2^{r-1}$, and its nonzero elements can have different
image lines inside $W$.

\paragraph{Precise unresolved geometric step.}
For a smooth simply connected compactification of the parameter space,
normal generation by divisorial inertia follows only after proving
that the finite quotient cover has no remaining divisorial branch and
applying purity. If the entire branch support were one irreducible
divisor, and its generic inertia were cyclic of order two, this would
give the single class required in the first lemma. More generally one
needs normal generation and conjugacy of all generating transvections,
not merely conjugacy of the generic inertia subgroups.

Neither the required branch-support statement nor cyclicity/conjugacy
has been proved for the remaining odd fixed-surface systems.
The monodromy reviewer reports that the degree-one relative determinant
already has the form $q(s)^2Q(s)$ with two cubic components, and that its
moving-kernel inertia is an elementary abelian group rather than a
single cyclic generator. The previously withdrawn higher-degree
alternating/Pfaffian model supplies no permissible higher-degree
factorization. Thus the existing odd-plus-even parabolic countermodel
is not excluded by the presently proved geometry.

The opened primary comparison was Brumer and Kramer, arXiv:1701.01890v1,
Theorem 2.1.1 and its reduction on printed p.3, and Proposition 2.3.4
with complete proof on printed p.5. Those are symplectic and arithmetic
statements; none is transferred to the present characteristic-two
unit-root family. The elementary criterion above is proved directly.
